\section{Recoil-point-cloud representation cross-check}
\label{sec:pet}

\paragraph{Status: diagnostic and method development, not a publication
uncertainty product.}
All results here are diagnostic and method development. No PET covariance,
historical recoil-only or full-event, is adopted as a publication uncertainty
product. The full-event central/statistical pairing was declined on 2026-08-20
and remains unresolved. Reconsideration requires estimator-equivalence and
coverage evidence; no branch was chosen
(App.~\ref{app:provenance}, \texttt{pet-pairing}).
% Moved to sec_summary.tex (Outlook), which is its single home (lane 5 R12); kept here as the
% source record, with the cross-reference corrected per RULINGS-F2 L3 #11 / F1 L5 C2:
% The publication uncertainty critical path ran through the
% adoption of a selection-complete scalar five-dimensional covariance, which has
% since happened under exception (\S\ref{sec:eavailw-adopted}), and never through this section;
% that path is not yet complete.
% Demoted from a rendered \TODO (lane 5 R12):
% TODO: Revise when the scalar-5D campaign (s5p, \texttt{OI-193}) closes. Adoption is not publication readiness: the s5e outcome records the publication objective as NOT READY.

\subsection{Inputs and output contract}

\paragraph{What goes in.}
\label{sec:pet-inputs}
PET is the point-cloud transformer used by OmniLearn~\cite{Mikuni:2024qsr},
vendored here inside OmniFold. Table~\ref{tab:pet-inputs} compares the historical
recoil-only cross-check with the later full-event extended-fiducial-phase-space
(extended-FPS, \S\ref{sec:fps}) candidate. Both classifiers use unbinned,
continuous event features without reporting-bin edges.

\begin{table}[H]
  \centering
  \footnotesize
  \setlength{\tabcolsep}{5pt}
  \renewcommand{\arraystretch}{1.08}
  \begin{tabular}{@{}>{\raggedright\arraybackslash}p{0.17\textwidth}
                      >{\raggedright\arraybackslash}p{0.375\textwidth}
                      >{\raggedright\arraybackslash}p{0.375\textwidth}@{}}
    \toprule
    object & historical recoil-only PET & later full-event extended-FPS PET \\
    \midrule
    reco/data tokens & Non-muon calorimetric clusters
      $(E_{\rm dep},\mathrm{pos}_{\rm view},z)$.  The transverse position is
      the coordinate in the cluster's X/U/V view. & The same generic
      calorimeter-cluster objects, extended to
      $(E_{\rm dep},\mathrm{pos}_{\rm view},z,\mathrm{view},t)$.  View is a
      detector-plane identifier, not an object-type label. \\
    \addlinespace[3pt]
    truth tokens & Final-state hadrons after removing the primary muon and
      neutrinos.  The dump stores $(E,p_x,p_y,p_z,\mathrm{PDG})$, but
      PDG is removed before classification
      (App.~\ref{app:provenance}, \texttt{pet-code}),
      so the classifier receives only the four-momentum. & Final-state-hadron
      tokens $(E,p_x,p_y,p_z,\mathrm{PDG},\theta,\cos\phi,\sin\phi)$; PDG is
      retained and azimuth is represented periodically. \\
    \addlinespace[3pt]
    event object & None.  Muon variables are used for event selection and for
      downstream reporting/binning, but are not classifier inputs. & Reco/data:
      $(\pT,\pPar)$, the reconstructed muon four-vector, $(\cos\phi,\sin\phi)$,
      $q/p$, the MINOS-match flag, and reconstructed vertex $(x,y,z)$.  Truth:
      muon $(\pT,\pPar)$.  The muon is a distinguished event-level input, not a
      cloud token, and conditions the transformer through feature-wise linear
      modulation (FiLM). \\
    \addlinespace[3pt]
    cardinality and grid & \multicolumn{2}{>{\raggedright\arraybackslash}p{0.75\textwidth}}{Clouds are
      energy-ranked and truncated or zero-padded to 12 tokens.  The extended
      $(\pT,\pPar)$ edges define domain retention, reporting, covariance and
      validation for the later estimator; they are never classifier features or
      training bins.} \\
    \bottomrule
  \end{tabular}
  \caption{Classifier input contract for the two PET representations used in this
  analysis.  ``Reco/data'' means that simulation at reconstruction level and
  measured data use the same observable schema; truth-only quantities are not
  manufactured on that leg.}
  \label{tab:pet-inputs}
\end{table}

Figure~\ref{fig:petdisplays} places representative reconstructed and truth
muons and recoil activity on a common energy--polar-angle plane. PET receives
variable-length sets of single-view energy deposits at reconstruction level
and particles with full four-momenta at truth level.
Figure~\ref{fig:petcardinality} shows each cloud's per-event cardinality.

The coverage-complete input in Figure~\ref{fig:petcardinality} has empty clouds
for only 174 of \num{32849103} truth events. GENIE residual-nucleus bookkeeping
particles are included in PET, giving mean truth cardinality $4.57$.
They also deplete the $k=1$ bin: a nominally single-hadron event, most often
quasielastic, usually contains a residual-nucleus pseudo-particle and enters
at $k=2$. Removing these pseudo-particles gives mean physical-hadron
multiplicity $4.43$ and a distribution rising smoothly to its $k=2$ maximum.
The twelve-constituent cap limits compute and memory costs, which grow with
cardinality. Only \SI{2.31}{\percent} of truth events reach it; the twelfth
constituent carries on average \SI{0.09}{\percent} of retained truth energy.

\begin{figure}[H]
  \centering
  \includegraphics[width=0.98\textwidth]{pet_event_displays_energy_angle}
  \caption{Four representative PET events on a common energy--polar-angle
  plane. Gray points mark reconstructed recoil clusters, colored stars truth
  final-state hadrons. Reconstructed (X) and truth (six-point star) muons nearly
  overlap. Muon momentum and direction come from the MINERvA--MINOS track and
  are stored as event-level $(\pT,\pPar)$ scalars. The illustrative cluster
  angle uses a nominal upstream vertex: the current reco-cloud dump stores
  $[E_{\rm dep},\mathrm{pos}_{\rm
  view},z]$ without the reconstructed vertex. The common projection is only
  a visual comparison. PET consumes raw single-view clusters at reco level
  and hadron four-vectors at truth level; both recoil-only inputs omit the
  muon (Table~\ref{tab:pet-inputs}).}
  \label{fig:petdisplays}
\end{figure}

\begin{figure}[H]
  \centering
  % _withremnant shows the truth cloud inputs to PET (including GENIE nuclear-remnant
  % pseudo-particles), which produces the depleted k=1 bin.
  % _real is the physical-hadron-only version.
  %   \includegraphics[width=0.72\textwidth]{pet_cardinality_real}
  \includegraphics[width=0.72\textwidth]{pet_cardinality_withremnant}
  \caption{POT-weighted per-event point-cloud cardinality (constituent count)
  for truth final-state hadrons and reconstructed clusters in MC and data.
  Reco-cluster multiplicity greatly exceeds truth hadron multiplicity because
  a single hadron deposits
  energy in many detector clusters. Data and MC reco agree closely
  (means \num{11.09} vs.\ \num{11.15}). Retaining only the twelve highest-energy
  constituents makes reco distributions pile up at the cap. Truth clouds are
  sparser (mean \num{4.57}, including GENIE nuclear-remnant bookkeeping
  particles) and rarely reach it. The text explains the depleted $k=1$ truth
  bin and the truncation choice.}
  \label{fig:petcardinality}
\end{figure}

\paragraph{What comes out.}
\label{sec:pet-output-contract}
The classifier score is not the scientific result. PET probabilities, implemented through raw logits, are intermediate
density-ratio estimates; PET emits no binned network prediction. OmniFold converts the final
truth-step logit to an unbinned push weight for each simulated truth event,
$w_{\rm push}=\exp(\mathrm{logit})$. The cross section follows by weighting the
stored truth inventory, histogramming the requested observable, and applying
standard flux, target and bin-volume normalization~\cite{Andreassen:2019cjw}.
Any supported function of the stored truth record or retained truth cloud can
then be projected without retraining. Projection cannot create configurations
outside generator support, recover constituents above the cloud cap, or
convert recoil-only weights to full-event weights.

\subsection{Shape cross-check against the scalar result}
\label{sec:pet-shape}
The first use of the recoil representation is a shape comparison against the production
scalar unfold, which tests the classifier family and the engineered-scalar versus
recoil-cloud input choice together.
% Here \texttt{cluster\_pos} is the measured coordinate in one X/U/V scintillator view.
Area-normalized against the frozen 4D GBDT result, the PET unfold
reproduces the scalar shape to a median \SIrange{2.3}{3.9}{\percent} per bin
($\pT$~\SI{3.9}{\percent}, $\pPar$~\SI{2.4}{\percent}, $\Eavail$~\SI{2.6}{\percent},
$q_{3}$~\SI{2.3}{\percent}; Fig.~\ref{fig:petgbdt}).
The remainder of this section takes the same
representation from this shape check to an absolutely-normalized cross section.

\begin{figure}[H]
  \centering
  \includegraphics[width=0.85\textwidth]{pet_vs_gbdt}
  \caption{Area-normalized 4D unfolded shapes from the
    recoil-only PET point-cloud transformer (solid)
    and production scalar GBDT (dashed), projected onto each axis. Median
    per-bin shape differences are \SIrange{2.3}{3.9}{\percent}. The only visible
    deviation, in the wide $q_{3}$ catch bin, is a binning artifact: the bin
    collects a large, structured high-$q_{3}$ tail and does not indicate a
    genuine shape disagreement. This is a shape-only comparison (PET on a
    2M-event subsample), not an independent physics result.}
  \label{fig:petgbdt}
\end{figure}

\subsection{From shape check to absolute cross section}
\label{sec:pet-abs}
PET trains on a GPU with a 2M-event subsample (2 OmniFold iterations).
Its step-2 truth reweight is applied to the \emph{full} 32.8M-event generator
sample, binned on the 4D grid and passed through standard extraction. The
extraction uses the GBDT run's reweight-independent completeness and the same
flux/POT/nucleon constants (App.~\ref{app:provenance}, \texttt{pet-code}).
The absolute path had two checks:

\begin{enumerate}[itemsep=2pt]
  \item Pushed weights over all 32.8M generator events are finite, with mean
        $1.028$.
  \item Ordinary closure used MC-reco-as-pseudo-data at only two iterations.
        It is a legacy self-consistency diagnostic, not a current validation;
        no current closure level is quoted
        (App.~\ref{app:provenance}, \texttt{pet-closure-legacy}). Pseudo-data
        and prior share the same simulation, so this test cannot expose
        omitted muon dependence at fixed recoil.
\end{enumerate}

\paragraph{No current PET--GBDT level of agreement is quoted.}
The only comparison used two-iteration training, preceded the 2026-08-01
full-event schema change and lacked background subtraction. It has not been
re-extracted under the pinned configuration. The legacy closure level and
higher-epoch retrain response are withdrawn on the same grounds
(App.~\ref{app:history}; App.~\ref{app:provenance}, \texttt{retracted-values}).

Ordinary closure alone cannot uniquely attribute a data-side gap to training
configuration. The recoil-only feature contract adds another difference:
at fixed recoil, the conditional muon distribution is inherited from the prior.
These structural limits survive the withdrawn numerical comparison.

\subsection{Statistical and training diagnostics of the recoil-only estimator}
\label{sec:pet-syst}
PET statistical replicas independently Poisson-fluctuate data and MC, retrain
at a fixed estimator seed, and use the same coherent MC draw for final truth
binning and completeness.

The replicas and nominal use scalar GBDT's background-subtracted
five-dimensional measured target, unlike the legacy comparison in
\S\ref{sec:pet-abs}. Data align exactly event by event; MC arrays are
byte-identical to the full-cloud input. Ordering, event-alignment and
finite-output gates passed for the nominal extraction, an identical-seed
GPU-floor repeat, a 20-member coherent-Poisson retraining ensemble, and a
12-member crossed subsample/estimator-seed ensemble. GPU nondeterminism is
negligible relative to sampling and training variation. The recoil-only
statistical block measures diagnostic dispersion from 20 coherent replicas;
a 100-replica ensemble has not been run. It is not an independently verified
covariance or the statistical uncertainty of any PET result, and is not paired
with a PET central value. More replicas would sharpen dispersion estimates
without supplying coverage evidence.

\emph{This block is not the full-event} $C_{\rm stat}$ \emph{of
App.~\ref{app:cstatlimit}} (App.~\ref{app:provenance}, \texttt{cstat-family}).
The constructions and populations differ: 20 recoil-only replicas versus
50 full-event members, with two resampled inventories versus three.
Neither object's measurements, disclaimers or rulings transfer to the other.

No uncertainty is quoted from the historical recoil-PET assembly of vertical,
retraining, ML, statistical and detector blocks (App.~\ref{app:history}).
% Cut here (lane 3 §9.3 #2); the reasons survive at App. H (app:history-petbudget), the
% "no current level" statements at sec:pet-abs:
% adding the separate frozen-map
% and retraining covariances omits their cross term, and the detector block is
% CV-support-limited (App.~\ref{app:history}).  This section therefore
% retains the recoil-only central extraction and ordinary closure as historical
% method-development diagnostics.  It quotes neither a current PET--GBDT agreement
% level nor a current closure level, and no recoil-only covariance is paired with
% a PET central value.

\subsection{Truth-cloud projection within the recoil cross-check}
\label{sec:pet-projection}

Each recoil-only PET truth event carries the generator-level hadron cloud
(top-12-by-energy hadrons with PDG codes, built by
\texttt{CVUniverse::GetTruthFSHadrons})
% truth cloud branch: \texttt{part\_gen}
alongside its learned step-2 push weight. Re-binning these weighted events
projects any observable of the stored 4-vectors without retraining or
re-unfolding. We demonstrate this \emph{reweight-then-project} capability on
three observables computed post-hoc from the stored cloud: available energy
$\Eavail$ (the exact operational definition in
\S\ref{sec:eavail-definition}), hadronic-system invariant mass
$W_{\text{had}}=\sqrt{(\sum E)^2-|\sum\vec p|^2}$, and proton/charged-pion
multiplicities, while quantifying the cloud's two limitations.

\paragraph{Truncation is essentially lossless.}
Cloud-recomputed and dump-time $\Eavail$ use the same formula and inputs,
so their residual isolates truncation to the 12 highest-energy hadrons.
Over all \num{32.8}~million truth events (every event now carries a cloud, below),
the residual has median \SI{0}{MeV}, mean \SI{\pcEavailMean}{MeV}, and
\SI{\pcEavailWithin}{\percent} of events within \SI{10}{MeV}
(Fig.~\ref{fig:pcvalid}). The deficit occurs entirely in the
\SI{\pcSaturated}{\percent} of events saturating the 12-particle cap
(median deficit \SI{\pcSatDeficit}{MeV} there, exactly zero below it). Energy ranking
removes the lowest-energy stored final-state hadrons, giving a small,
one-sided, bounded effect. The 12-particle cap therefore does not limit $\Eavail$-scale
observables and does not need raising.

\paragraph{Coverage: closing the acceptance gap.}
The truth-only ``miss'' events appended for the unfold (\SI{\pcMiss}{\percent} of
the truth sample) carry a cloud: the truth final-state hadrons are carried through the
truth-denominator loop, filling the truth cloud for the appended miss
rows (the reco cloud correctly stays empty for a miss;
% branches: truth cloud \texttt{part\_gen}, reco cloud \texttt{part\_reco}
App.~\ref{app:history}). Coverage is complete:
\SI{\pcHasCloud}{\percent} of truth events carry clouds, leaving only
\pcEmptyCloud{} residual empties out of \num{32.8}~million. Cloud-recomputed
$d\sigma/d\Eavail$ reproduces the full unfold to
$\lesssim\SI{\pcCloudFull}{\percent}$ in every $\Eavail$ bin; the projection
is no longer acceptance-limited.

Three $\Eavail$ projections overlay in the centre panel of
Fig.~\ref{fig:pcxsec}: push-weighted cloud and stored $\Eavail$ on the
has-cloud subset, and stored $\Eavail$ over the full unfold. Their largest
relative per-bin spread is
\SI{\pcCloudFull}{\percent}. No reco-efficiency division is applied because
unfolded counts are already acceptance-corrected
(App.~\ref{app:history}; App.~\ref{app:provenance}, \texttt{pet-projection-figure}).
% Summary keys: \texttt{pet\_cloud\_hascloud}, \texttt{pet\_stored\_hascloud}, \texttt{pet\_stored\_full}.

\paragraph{Exploratory projection yields new hadronic diagnostics.}
$W_{\text{had}}$ measures the true hadron-system invariant mass. The
leptonic ``experimenter's'' $W$ stored at dump time is reconstructed from muon kinematics assuming
a single struck nucleon at rest; $W_{\text{had}}$ is not its lossy copy.
Their median difference is \SI{\pcWoffset}{GeV}. The 2D correlation
(Fig.~\ref{fig:pcxsec}, left) has a sharp locus at
$W_{\text{had}}\approx\SI{0.94}{GeV}$ for single-baryon final states: their
hadronic mass equals the nucleon rest mass irrespective of leptonic $W$.
Fermi-smeared diagonal bands rise by $\sim$one nucleon mass per extra
final-state baryon (multi-nucleon knockout, 2p2h). The projection thus supplies
a physically interpretable hadron-cloud diagnostic from the recoil-only unfold.

\begin{figure}[H]
  \centering
  \includegraphics[width=0.90\textwidth]{pet_cloud_projection_validation}
  \caption{Cloud-recomputed $\Eavail$ minus its dump-time value over all
  \num{32.8}~million truth events. The residual is zero to
  \SI{\pcEavailWithin}{\percent} within \SI{10}{MeV}; its small negative tail
  occurs only in the \SI{\pcSaturated}{\percent} reaching the 12-hadron cap.
  Right: ordinary event counts versus pre-truncation multiplicity from the
  variable-length ROOT branch. Multiplicities 1--12 are exact; sparse tails
  are grouped into labeled 13--14, 15--17, 18--24, and $\geq25$ categories.
  The dashed boundary separates fully stored events from categories retaining
  only the 12 highest-energy hadrons. Cap-reaching events have median
  $\Eavail$ deficit \SI{\pcSatDeficit}{MeV}. Keeping the 12 highest-energy hadrons is therefore
  essentially lossless for $\Eavail$-scale observables.}
  \label{fig:pcvalid}
\end{figure}

\begin{figure}[H]
  \centering
  \includegraphics[width=0.98\textwidth]{pet_cloud_projection_xsec}
  \caption{Exploratory recoil-only reweight-then-project demonstration.
  \emph{Left:} hadronic-system invariant mass $W_{\text{had}}$ versus stored
  leptonic $W$. Their definitions differ (median \SI{\pcWoffset}{GeV}); the
  QE locus is at \SI{0.94}{GeV}, with multi-baryon bands above $y=x$.
  \emph{Center:} projected $d\sigma/d\Eavail$. Push-weighted cloud $\Eavail$,
  stored $\Eavail$ on the has-cloud subset, and stored $\Eavail$ over the full
  unfold agree to $\lesssim\SI{\pcCloudFull}{\percent}$ per bin because
  essentially every truth event carries a cloud after the coverage fix.
  The different GBDT 5D estimator is not expected to coincide.
  \emph{Right:} projected $d\sigma/d W_{\text{had}}$, whose shape differs
  from the leptonic-$W$ reference as expected.}
  \label{fig:pcxsec}
\end{figure}

\paragraph{The projection and extrapolation limits.}
\label{sec:pet-projection-next}
Reweighting can only redistribute weight among truth configurations the
generator already produced; it cannot populate phase space the prior never
sampled --- the same extrapolation limit flagged for the extended-fiducial
study (\S\ref{sec:fps}). Hadron-cloud diagnostics can be projected within that
stored spectrum. They cannot convert recoil-derived weights to full-event
weights or remove the prior conditional for omitted muon information.
The stored recoil-only result cannot acquire the later full-event estimator's
muon dependence: those are separate estimators with the
distinct interfaces in Table~\ref{tab:pet-inputs}.
Re-projection also cannot recover rare very-high-multiplicity constituents
above the 12-particle cap, which the validation finds unnecessary at
$\Eavail$ scale. Details and reproduction are in App.~\ref{app:provenance}
(\texttt{pet-projection-figure}).

\subsection{Extended population and the full-event candidate}
\label{sec:pet-full-event-status}

\paragraph{Recoil-PET feasibility study on the extended population.}
\label{sec:fps-pet}
The feasibility study on the enlarged population and reporting grid of
\S\ref{sec:fps} remained recoil-only. Step~1 used non-muon reco clusters;
step~2 used truth hadrons with the muon removed. Muon coordinates defined
and binned the extension but were absent from classifier inputs. The
populated-bin counts, prior-transfer envelope and retraining floor are
recoil-estimator feasibility diagnostics, not a final full-event measurement;
they are not used to argue that low-efficiency cells are data-constrained.

A full-event estimator on this population would carry the requirements set out
above, and additionally consume the canonical extended-fiducial event loops
and $15\times19$ muon grid while preserving the enlarged truth denominator, native
misses and unchanged reconstructed selection. Additional features can use more observed
information, but cannot create detector information in zero-efficiency cells.

\paragraph{Full-event artifacts and present boundary.}
Within its recoil-only contract, the study demonstrates end-to-end paired reco/truth cloud
construction, GPU training, full-sample reweighting, absolute cross-section
extraction, coherent bootstrap replicas and the response studies above.
It does not establish simultaneous muon-plus-hadron unfolding, because the muon
was absent from classifier inputs.

Later full-event extended-FPS work has a nominal artifact and a 50-member
$C_{\rm stat}$ artifact. The latter has one builder, is not independently
verified or adopted, and was declined for pairing with the nominal
(status paragraph above).
% Joseph declined pairing it with the nominal on 2026-08-20.
% Cut as a restatement of the status paragraph (lane 5 R8): "and requires
% estimator-equivalence plus coverage evidence before reconsideration" and "and no PET
% covariance is adopted as a publication uncertainty". Gate 6 and the five receipt
% prohibitions moved to App. E (pet-gate6).
The ruling neither invalidates the measured dispersion nor establishes a
mechanism. A spatially coherent anomaly remains with unvalidated coverage.
No $C_{\rm ML}$ or full PET total covariance exists, and starting a PET
ML-variance component is not authorized
(App.~\ref{app:provenance}, \texttt{pet-gate6}). App.~\ref{app:cstatlimit}
records the construction and limitations. None of these artifacts changes PET's diagnostic and method-development status.

\subsection{Related representations and estimator-design outlook}
\label{sec:pet-gregor}
Gregor Krzmanc's \texttt{minerva-ml} uses the same broad MINERvA ME~FHC
simulation for supervised available-energy regression and charged-current-pion
classification with an OmniLearned/PET2 backbone~\cite{Krzmanc:2026transfer}.
At the pinned commit, typed muon, photon, prong and calorimeter-blob tokens
use collider-style $(\eta,\phi,\log p_T,\log E)$ coordinates, learned type
embeddings, $dE/dx$, position and timing. Prongs or blobs exceeding their
separate caps become aggregate overflow tokens.\footnote{Implementation comparison:
\href{https://github.com/gregorkrz/minerva-ml/blob/af5d92ed2b3b448a09b6b7cf6b4f179e5757b4ed/DATASET.md}{\texttt{minerva-ml DATASET.md}}
and
\href{https://github.com/gregorkrz/minerva-ml/blob/af5d92ed2b3b448a09b6b7cf6b4f179e5757b4ed/MODELS.md}{\texttt{MODELS.md}},
commit \texttt{af5d92ed}.}
Our reco leg uses generic detector-geometry calorimeter clusters without
reconstructed-object type embeddings and discards tokens beyond the common
top-12-by-energy cap. Our task estimates OmniFold density ratios against real
data with explicit backgrounds and native misses, rather than supervised
regression or classification on simulation. Typed reconstructed objects and
aggregate overflow tokens remain prospective representation improvements;
the present estimator uses neither.

\paragraph{Ensembling as an estimator-design choice.}
Neural likelihood-ratio studies find that Step~1 model averaging before
pushing weights to Step~2 can improve unfolding bias and variance
\cite{ToralesAcosta:2025stabilizing}. Averaging within iterations and averaging
final weights define different estimators. Future comparisons must specify
what is averaged, where in the iteration, and the nominal/replica pairing.
The same study finds that pretraining can reduce
variance while increasing bias, so seed dispersion alone is insufficient to
rank designs. Closure bias and uncertainty calibration must also be assessed.
These are method-development directions, not explanations of the existing
full-event anomaly, evidence for its declined pairing, or an adopted PET
uncertainty construction.
